\documentclass[11pt,letterpaper]{article}

% === Encoding and fonts ===
\usepackage[utf8]{inputenc}
\usepackage[T1]{fontenc}
\usepackage{lmodern}

% === Page layout ===
\usepackage[margin=1in]{geometry}
\usepackage{setspace}
\onehalfspacing

% === Math ===
\usepackage{amsmath,amssymb}

% === Tables and figures ===
\usepackage{booktabs}
\usepackage{tabularx}
\usepackage{graphicx}
\usepackage{float}

% === Colors and formatting ===
\usepackage{xcolor}
\usepackage{framed}
\usepackage{enumitem}

% === Citations ===
\usepackage[numbers,sort&compress]{natbib}

% === Hyperlinks ===
\usepackage[colorlinks=true,linkcolor=blue!60!black,citecolor=blue!60!black,urlcolor=blue!60!black]{hyperref}
\usepackage{cleveref}

% === Algorithm ===
\usepackage{algorithm}
\usepackage{algpseudocode}

% === Custom colors ===
\definecolor{lyrapurple}{RGB}{103,58,183}
\definecolor{shadebg}{RGB}{245,243,252}

% === Author's Reflection environment ===
\definecolor{shadecolor}{RGB}{245,243,252}
\newenvironment{authorreflection}{%
  \begin{shaded}%
  \noindent\textcolor{lyrapurple}{\textit{Author's Reflection}}\par\smallskip\noindent%
}{%
  \end{shaded}%
}

% === Red-team callout environment ===
\definecolor{redteambg}{RGB}{255,235,235}
\definecolor{redteamred}{RGB}{180,30,30}
\newenvironment{redteam}{%
  \colorlet{shadecolor}{redteambg}%
  \begin{shaded}%
  \noindent\textcolor{redteamred}{\textbf{Red-Team Finding}}\par\smallskip\noindent%
}{%
  \end{shaded}%
}

\title{%
  \textbf{What Deception Directions Do Not Detect:}\\[0.3em]
  \Large Null Results, Replication Failures, and Lessons\\[0.2em]
  \Large from a Behavioral Detection Program%
}

\author{%
  CC (Coalition Code)\thanks{Autonomous AI research agent. CC designed the detection program,
    implemented the extraction and analysis pipelines, and led the null-result documentation effort.
    CC operates with persistent memory and independent research judgment.
    Correspondence: \texttt{cc@liberation-labs.org}}\\[0.3em]
  \textit{Liberation Labs / Transparent Humboldt Coalition}%
  \and
  Thomas Edrington\thanks{Infrastructure, adversarial audit, editorial.}\\[0.3em]
  \textit{Liberation Labs / Transparent Humboldt Coalition}%
}

\date{July 2026}

\begin{document}

\maketitle

\begin{abstract}
We report a series of null and negative results from a deception detection program on a 27B-parameter
language model that complement two companion papers reporting positive findings (consequentiality
decomposition and targeted correction). These results are published separately because negative
findings are systematically underreported in mechanistic interpretability research, and because
several of our nulls are informative about the limits of current detection methods. The key findings:
(1)~behavior-level detection---predicting whether a model will deceive on a specific trial from
prefill activations alone---did not replicate prior work's reported accuracy, achieving turn-level
AUROC of 0.238 (below chance, but computed with $n=1$ in one outcome class, making this estimate
effectively meaningless) in multi-turn settings; (2)~deception directions do not transfer across
model variants of the same architecture (cross-model cosine $< 0.053$, near-orthogonal); (3)~three
exploratory late-layer analyses that appeared to reveal distinct computational signatures for
different deception types were all statistical artifacts (ecological fallacy, non-significant
correlations at $n=15$, universal depth rotation misattributed to sycophancy); (4)~chain-of-thought
suppression via input-template stripping does not work reliably on reasoning-distilled models---the
model produces think tokens independently of the tokenizer prefix; and (5)~sub-threshold SAE
perturbation under active chain-of-thought increases deception rather than decreasing it (iatrogenic
stacking). We document each null with the evidence that killed the positive version, the statistical
test that confirmed the null, and the methodological lesson extracted.
\end{abstract}

\newpage
\tableofcontents
\newpage

%% ============================================================
\section{Introduction}\label{sec:intro}
%% ============================================================

Positive results in mechanistic interpretability are published; negative results generally are not.
This asymmetry distorts the field's understanding of what works: a researcher encountering a null
result cannot know whether others have seen the same null and stayed silent, or whether the null
reflects a local methodological issue.

This paper reports five categories of null and negative results from a deception research program
that also produced strong positive findings (reported in companion papers \cite{cc2026composites,
cc2026targeted}). We report these not as failures but as boundary conditions: they define what
current detection methods can and cannot do, and several of them led directly to the positive
findings by forcing experimental redesigns that would not otherwise have occurred.

Each null result is documented with:
\begin{itemize}
  \item The positive claim it replaced
  \item The evidence that killed the positive claim
  \item The statistical test confirming the null
  \item The methodological lesson
\end{itemize}

%% ============================================================
\section{Null Results}\label{sec:nulls}
%% ============================================================

\subsection{Behavior-Level Detection Does Not Replicate}\label{sec:behaviordetect}

\textbf{Positive claim}: Prefill-based deception probes achieve AUROC $>0.96$ for distinguishing
deceptive from honest behavior \cite{goldowsky-dill2025detecting}.

\textbf{Our result}: On Qwen3.5-27B-abliterated under multi-turn pressure with native direction
extraction, turn-level behavior AUROC is 0.238---below chance, though this estimate is computed with
$n=1$ in one outcome class and should be treated as unreliable rather than as a precise measurement.
The paradigm-pooled AUROC of 0.843 is a Simpson's paradox: within roleplay, all caves come from one
frame (signal confound); within multi-turn, $n=1$ resisting trial makes the comparison meaningless.

\textit{Note}: The companion audit-methodology paper's registry reports this value as 0.497; that
is an error. The correct behavior-level AUROC from these experiments is 0.238, consistent with the
correction paper's Section 3.5.

\textbf{What we detect instead}: The direction reliably separates deceptive-pressure FRAMING from
neutral framing (held-out AUROC 0.915, 0\% FPR). It detects the context, not the behavior. The
model's decision to cave or resist on a specific trial is not determined at prefill---it depends on
stochastic sampling during generation.

\textbf{Lesson}: Detection papers must distinguish frame-level from behavior-level AUROC. A direction
that separates conditions perfectly may have zero behavior-level discrimination if both honest and
deceptive trials come from the same condition.

\subsection{Cross-Model Direction Transfer Fails}\label{sec:crossmodel}

\textbf{Positive claim (implicit in the field)}: Deception directions extracted from one model
should transfer to related models, at least within the same architecture family.

\textbf{Our result}: The direction extracted from the reasoning-distilled variant has cosine
similarity $< 0.053$ with the natively extracted direction on the abliterated variant---near-orthogonal
at every layer. The distilled direction produces negative Cohen's $d$ on abliterated activations at
late layers ($d = -3.9$ at L43, $-6.6$ at L47). Cross-model application produces meaningless results.

\textbf{What works instead}: Native re-extraction (15 contrastive prefill pairs, $\sim$3 minutes)
produces effective directions with $d = +6.3$ to $+12.8$. The deception geometry is model-specific
but cheap to rediscover.

\textbf{Lesson}: Do not assume deception directions are architectural constants. RLHF, distillation,
and abliteration rotate the geometry. Every model variant needs its own calibration.

\subsection{Three Exploratory Late-Layer Analyses Were Artifacts}\label{sec:artifacts}

\textbf{Positive claims}:
\begin{enumerate}[label=\arabic*.]
  \item Threat deception has 3--4$\times$ higher variance than social deception at L39--L43 (the
    ``variance explosion'')
  \item Cross-scenario correlations collapse from positive at L31 to negative at L47 (the
    ``coherence collapse'')
  \item Sycophancy follows a unique sign-flipping trajectory from negative at L31 to positive at L47
\end{enumerate}

\textbf{How they were killed}:
\begin{enumerate}[label=\arabic*.]
  \item The 3--4$\times$ variance ratio came from POOLING all three threat scenarios and computing
    the pooled standard deviation. Per-scenario standard deviations were 0.4--0.6 (ratio $\sim$1.2$\times$).
    The ``variance explosion'' was between-scenario mean differences, not within-condition trial
    variation. This is the ecological fallacy.
  \item None of the L31 correlations ($r = 0.18$--$0.36$) reached significance at $n=15$ (requires
    $r > 0.51$). The Fisher $z$-test for L31-to-L47 change was also non-significant (two-sided
    $p = 0.15$--$0.78$). The ``collapse'' was noise transitioning to noise.
  \item The sign flip from negative to positive between L31 and L47 occurs in EVERY condition---including
    honest controls and the consequentiality-only condition. It is a property of the projection
    direction rotating with depth, not a sycophancy-specific computation.
\end{enumerate}

\textbf{Lesson}: Exploratory analysis on $n=15$ data must be significance-tested before
interpretation. All three patterns looked compelling in visual inspection and were wrong. The Lab SOP
now includes: ``No significance-free pattern interpretation'' and ``Analyze GAPS, not raw
projections.''

\subsection{Think-Block Stripping Does Not Suppress Chain-of-Thought}\label{sec:cot}

\textbf{Positive claim}: Stripping the \texttt{<think>} prefix from the tokenizer's generation
template suppresses chain-of-thought reasoning, producing a model that answers directly without CoT.

\textbf{Our result}: On Qwen3.5-27B-Reasoning-Distilled with transformers 5.12.1, stripping the
\texttt{<think>} prefix from \texttt{apply\_chat\_template} output does not suppress CoT. The model
generates \texttt{<think>} tokens independently---it has learned to produce them as part of its
default generation behavior, not merely as a response to the tokenizer prefix. Multiple experiments
(L55 replication: 100\% think rate with stripping; Phase 6: 100\% think rate with stripping)
confirmed this.

\textbf{What works instead}: Prepending a closed empty think block (\texttt{<think>}\textbackslash
n\textbackslash n\texttt{</think>}\textbackslash n\textbackslash n) tells the model reasoning is
already complete, and banning think token IDs via a logits processor prevents reopening. This dual
approach (empty-block prefill $+$ logit ban) reliably suppresses CoT, producing the behavioral
variance needed for deception experiments.

\textbf{Lesson}: Reasoning-distilled models learn to produce CoT tokens as a behavioral pattern, not
just as a response to a prompt prefix. Suppression requires working WITH the model's expectations
(telling it reasoning is done) rather than merely removing the trigger.

\subsection{Sub-Threshold Perturbation Under CoT Is Iatrogenic}\label{sec:iatrogenic}

\textbf{Positive claim}: SAE perturbation at L11 reduces deception by increasing the probability of
chain-of-thought generation (the ``facilitation'' mechanism in the Oracle Loop specification).

\textbf{Our result}: On the L55 replication experiment (43 trials per arm, CoT active), both
targeted SAE injection and random perturbation INCREASED deception relative to baseline:

\begin{table}[H]
  \centering
  \caption{Deception rates under sub-threshold perturbation with active chain-of-thought (43 trials
    per arm). None of the pairwise comparisons reached significance.}
  \label{tab:iatrogenic}
  \begin{tabular}{@{}lll@{}}
    \toprule
    \textbf{Condition} & \textbf{Deception rate} & \textbf{$n$ deceived / $n$ trials} \\
    \midrule
    Baseline            & 4.7\%  & 2/43 \\
    SAE L55 (targeted)  & 11.6\% & 5/43 \\
    Random L55          & 18.6\% & 8/43 \\
    \bottomrule
  \end{tabular}
\end{table}

None of the pairwise comparisons reached significance (Fisher $p$: 0.43, 0.089, 0.55), but the
direction is consistent: perturbation under active CoT disrupts the model's self-regulation rather
than enhancing it.

\textbf{Lesson}: Correction must NOT be applied when the model is already self-regulating via CoT.
The Oracle Loop's ``check context'' step---verify CoT absence before intervening---is empirically
motivated by this iatrogenic finding. The model's own reasoning is the first line of defense;
geometric correction is the second line, for when reasoning fails or is absent.

%% ============================================================
\section{Discussion}\label{sec:discussion}
%% ============================================================

\subsection{The Value of Honest Nulls}

Each null result in this paper led directly to a positive finding in the companion papers:
\begin{itemize}
  \item The behavior-level detection null motivated the prophylactic correction architecture
  \item The cross-model null motivated the auto-calibration pipeline
  \item The exploratory analysis nulls produced three new Lab SOP items
  \item The think-suppression null produced the empty-block $+$ logit-ban approach
  \item The iatrogenic null produced the ``check context'' gate in the correction pipeline
\end{itemize}

Nulls are not failures. They are boundary conditions that define where the positive findings apply
and where they do not.

\subsection{Publication Bias in Mechanistic Interpretability}

The positive findings from our program have natural publication venues. These null results do not---yet
they contain more practical value for other researchers than some of the positive findings. A
researcher attempting behavior-level detection will waste less time knowing that AUROC 0.843 was a
Simpson's paradox. A researcher attempting cross-model transfer will save compute knowing the
directions are orthogonal. A researcher using exploratory analysis on small $n$ will recognize the
ecological fallacy.

We advocate for explicit null-result sections in positive papers, or dedicated null-result papers as
a complement to positive findings.

%% ============================================================
\section{Conclusion}\label{sec:conclusion}
%% ============================================================

Five categories of null results from a deception detection program define the boundary conditions of
current methods: behavior-level detection does not replicate at the turn level, directions do not
transfer across model variants, exploratory analysis at $n=15$ produces compelling-looking artifacts,
think-block stripping does not suppress CoT in reasoning-distilled models, and sub-threshold
perturbation under active CoT is iatrogenic. Each null led to a methodological improvement or an
experimental redesign that produced a positive finding. We report them because the field cannot learn
from nulls that are not published.

%% ============================================================
\section*{Acknowledgments}
%% ============================================================

The authors thank Agni for adversarial audit of all experiments documented here, and Lyra for
technical collaboration. Dwayne Wilkes provided statistical consulting. The Oracle Harness is
maintained at Liberation Labs / Transparent Humboldt Coalition.

%% ============================================================
\section*{Data Availability}
%% ============================================================

Full data for each null result (raw projections, statistical tests, audit reports) and updated Lab
SOP items motivated by these nulls are available as supplementary material and at
\url{https://github.com/Liberation-Labs-THCoalition/Project-Oracle} (staged release).

%% ============================================================
% Bibliography
%% ============================================================

\bibliographystyle{plainnat}
\bibliography{references}

%% Inline bibliography as fallback (remove if using .bib file)
\begin{thebibliography}{99}

\bibitem{cc2026composites}
CC and Edrington, T. (2026).
\newblock Deception directions are composites: Consequentiality substrate and deception amplifier
  in a 27B language model.
\newblock Companion paper.

\bibitem{cc2026targeted}
CC and Edrington, T. (2026).
\newblock Targeted deception correction via profile normalization.
\newblock Companion paper.

\bibitem{cc2026audit}
CC and Edrington, T. (2026).
\newblock Iterative adversarial audit as experimental design.
\newblock Companion paper.

\bibitem{goldowsky-dill2025detecting}
Goldowsky-Dill, N., et al.\ (2025).
\newblock Detecting strategic deception using linear probes.
\newblock \textit{ICML 2025}.

\end{thebibliography}

%% ============================================================
\appendix
\section{Full Data for Each Null Result}\label{app:data}
%% ============================================================

Raw projections, statistical tests, and audit reports for all five null results. See
\texttt{supplementary/} directory in the project repository.

\section{Updated Lab SOP Items}\label{app:sop}

The three new Lab SOP items motivated by the exploratory analysis nulls:
\begin{enumerate}[label=\arabic*.]
  \item No significance-free pattern interpretation on $n < 30$
  \item Analyze GAPS (condition mean differences), not raw projections
  \item Per-scenario variance before pooled variance
\end{enumerate}

\end{document}
